\chapter{System Requirements and Problem Analysis}

The Mini Payment Gateway is positioned as a pilot-ready payment gateway rather
than a full production payment platform. Its purpose is to support merchant
onboarding, VietQR-based payment creation, asynchronous provider result
processing, full refunds, webhook delivery, operational investigation, and
merchant self-visibility through a separate dashboard experience.

\section{Problem Context}

The gateway solves two connected problems. First, merchants need a small but
reliable payment interface for creating QR payments, receiving asynchronous
provider outcomes, querying transaction status, and requesting full refunds.
Second, internal operators need visibility and recovery tools because payment
systems are asynchronous by nature: callbacks can arrive late, webhook delivery
can fail, and merchant readiness must be enforced before money-flow APIs are
usable.

The VietQR pilot adds one more practical problem: a merchant payment is not
usable until the gateway knows which receiving bank account should be encoded
into the QR payload. For this reason, QR receiving accounts are controlled by
internal operations rather than by merchant self-service. This keeps pilot
configuration auditable and makes payment creation fail early when a merchant
is active but not yet configured for scannable VietQR payments.

The dashboard work revealed an architectural constraint that strongly shaped the
project. A merchant-facing dashboard cannot safely be implemented as frontend
screens alone. It requires a separate merchant-facing access model, merchant-
scoped read capability, and internal provisioning support. Without these
boundaries, the system would either misuse merchant integration credentials for
human login or risk cross-merchant data exposure.

\section{Implemented Scope And Non-Goals}

\begin{table}[!htbp]
\centering
\caption{Implemented system scope}
\begin{tabularx}{\textwidth}{L{0.24\textwidth}Y}
\toprule
\textbf{Area} & \textbf{Implemented capability} \\
\midrule
Merchant lifecycle & Create merchants, submit and review onboarding evidence, create or rotate credentials, activate, suspend, and disable merchants. \\
\bodyrowrule
Payments & Create dynamic VietQR payments, enforce semantic idempotency, return QR reference, EMV-compatible QR payload, and base64 PNG image, query by transaction id or order id, and process provider payment result callbacks. \\
\bodyrowrule
Refunds & Create full refunds, enforce refund rules, query refund state, and process provider refund result callbacks. \\
\bodyrowrule
Webhooks & Create durable signed outbound webhook events, store delivery attempts, apply retry policy, and support internal manual retry. \\
\bodyrowrule
Ops Dashboard & Internal operator interface for auth, RBAC, merchant workflows, VietQR receiving account management, evidence inspection, reconciliation, audit, internal users, and portal-user provisioning. \\
\bodyrowrule
Merchant Dashboard & Merchant-scoped read-only interface for overview, analytics, payments, refunds, webhooks, profile, credentials, and password change. \\
\bodyrowrule
Testing & Automated business-rule checks, end-to-end transaction scenarios, dashboard validation, and browser-based usability review. \\
\bottomrule
\end{tabularx}
\end{table}

The following items are intentionally out of scope for the implemented pilot:

\begin{itemize}
  \item settlement ledger and accounting workflows;
  \item disputes or chargebacks;
  \item partial refunds;
  \item multi-provider routing;
  \item merchant self-service onboarding;
  \item merchant self-service QR receiving account management;
  \item production hosted checkout pages (the repository includes only a local
  demo merchant integration for presentation and verification);
  \item merchant profile editing from the Merchant Dashboard;
  \item real bank/provider settlement integration;
  \item production-grade notification, TLS, and observability hardening.
\end{itemize}

\section{Functional Requirements}

\begin{longtable}{L{0.24\textwidth}L{0.30\textwidth}L{0.36\textwidth}}
\caption{Functional requirements}\\
\toprule
\textbf{Requirement group} & \textbf{Rule} & \textbf{Implemented decision} \\
\midrule
\endfirsthead
\toprule
\textbf{Requirement group} & \textbf{Rule} & \textbf{Implemented decision} \\
\midrule
\endhead
Merchant lifecycle & Merchant statuses are pending review, active, rejected, suspended, and disabled. & Merchant readiness review and merchant operating status are tracked separately so approval outcome and live operating state do not become confused. \\
\bodyrowrule
Payment creation & One active payment is allowed per merchant and order. & A pending payment is reused only when the duplicate request is semantically identical; a completed successful payment blocks a new attempt. \\
\bodyrowrule
QR payment readiness & An active merchant must have one active VietQR receiving account before creating a payment. & Ops configures the receiving bank, account number, account name, and template. Payment creation returns \code{qr_reference}, \code{qr_content}, and \code{qr_image_base64}. \\
\bodyrowrule
Pilot QR constraints & The pilot supports VND whole-number payments only. & The transfer purpose is short, ASCII-safe, and derived from the gateway QR reference so it can be matched against callback evidence. \\
\bodyrowrule
Payment states & Persisted states are pending, successful, failed, and expired. & External outcome evidence finalizes pending payments, while late or mismatched evidence is routed to reconciliation. \\
\bodyrowrule
Refunds & Only full refunds are implemented; refund eligibility expires after a limited business window. & Refund records move through pending, refunded, and refund failed states. Partial refunds are rejected by design. \\
\bodyrowrule
Webhooks & Successful merchant acknowledgement marks delivery success; failed delivery activates controlled retry. & Delivery attempts are preserved as evidence; exhausted events become failed and can later be retried by internal operators when eligible. \\
\bodyrowrule
Provider callbacks & External payment and refund callbacks must be authenticated. & The simulator/provider callback surface requires provider id, timestamp, and HMAC signature headers before applying state transitions. \\
\bodyrowrule
Dashboards & Ops and Merchant users require different surfaces. & Ops Dashboard is internal and mutable; Merchant Dashboard is read-only except password change. \\
\bottomrule
\end{longtable}

\section{Security And Data Requirements}

Each access boundary uses a matching authentication model. Merchant backends use
signed integration credentials. Internal users authenticate through an internal
session model. Merchant dashboard users authenticate through a separate
merchant-scoped session model. This separation prevents human-facing access from
reusing merchant integration secrets and keeps merchant visibility within the
correct tenant boundary.

The system also enforces persistence rules where practical:

\begin{itemize}
  \item one active credential per merchant;
  \item one active VietQR receiving account per merchant and provider;
  \item one pending payment per merchant and order;
  \item one logical refund per merchant refund reference;
  \item at most one successful full refund per payment;
  \item positive monetary values and non-negative attempt counters;
  \item merchant-scoped analytics indexes for payment, refund, and webhook
  records.
\end{itemize}

\section{Operational Requirements}

The gateway must remain supportable by internal operators. Every important
manual action carries an audit reason, and audit rows capture the event code,
entity identity, actor type, actor id, before and after snapshots, and
timestamp. Sensitive fields such as raw credential secrets are masked in
responses and audit snapshots.

The merchant-facing dashboard must help merchants inspect their own data without
seeing internal reconciliation workflows, raw credential secrets, or data owned
by other merchants. This requirement drives the separate merchant-facing access
boundary and the rule that post-login dashboard reads never rely on a
client-supplied merchant identifier as the source of scope.
